\section{An adaptive separator tree in polynomial space}
\label{sec:adaptive-separator-new}

The same second-moment idea gives a separate improvement that retains
polynomial space. It changes only the order of the exact group-factor sum;
it needs no matrix multiplication. We use the existing streamed separator
algorithm and its $r=16$ precision and packing schedule. Its earlier
coefficient is $55.709478318910\ldots$.

\begin{theorem}[Adaptive streamed coefficient]
\label{thm:adaptive-separator-new}
Under the same normalized contraction interface and inherited algebraic
identities, there is a polynomial-space algorithm with time
\[
 (s+\mathcal L)^{O(1)}
 2^{(C_{\rm ps}+o(1))s/\sqrt{\log_2s}},\qquad
 C_{\rm ps}=\frac{5\pi}{\ln2}\sqrt{\frac{145}{42}}
 \left(1+\frac1{3\sqrt3}\right)
 =50.210434836331\ldots .
\]
In particular, $50.22$ is a safe eventual coefficient.
\end{theorem}

\paragraph{The exact local sum.}
Fix a pair-budget array for sixteen groups. Each directed pair $e$ has
an exact-weight index list of size $N_e$, and write $h_e=\log_2N_e$.
Empty lists discard the array. The imported entropy bound, in its
Euclidean form, gives
\[
 \sum_e h_e^2\le (1+o(1))\frac{K^2}{16^2},\qquad
 K=\frac{\pi\sqrt2}{\ln2}\frac{M}{\sqrt u}.
\]
For each unordered pair let $w_{ab}=h_{ab}+h_{ba}$, and put
$q=\sum_{a<b}w_{ab}^2$ and $W=\sum_{a<b}w_{ab}$. Thus
$q\le(1+o(1))K^2/128$.

Take the full balanced binary tree on sixteen leaves. At a node, condition
on all directed indices crossing its two child blocks, recurse independently
on the two sides, multiply the answers, and sum. For a leaf $\ell$, let
$E_\ell$ be the union of the cut-pair sets on its root-to-leaf path, viewed
as unordered pairs, and let
\[
 L_\ell=\sum_{e\in E_\ell}w_e.
\]
The number of paid group-factor evaluations is at most
$\sum_\ell2^{L_\ell}\le16\,2^{\max_\ell L_\ell}$, as in the inherited
streamed separator identity. Each group-factor evaluation is an actual
recursive child contraction. No cache or free oracle is used.

\begin{lemma}[Adaptive tree moment bound]
\label{lem:adaptive-tree-moment-new}
For every nonnegative weighting of the $120$ unordered pairs, some
relabeling of the balanced sixteen-leaf tree satisfies
\[
 \max_\ell L_\ell
 \le \sqrt{\frac{725}{1344}}\,K+o(K).
\]
\end{lemma}

\begin{proof}
Uniformly permute the sixteen group labels over the leaves. A leaf path
contains $64+16+4+1=85$ unordered cut pairs. Write
\[
 m=\frac1{16}\sum_\ell L_\ell,\qquad
 V=\frac1{16}\sum_\ell(L_\ell-m)^2.
\]
Linearity of expectation gives $\mathbb E m=(85/120)W=(17/24)W$.
For any sixteen real numbers,
\[
 \max_\ell L_\ell\le m+\sqrt{15}\sqrt V.
\]
Indeed, if the largest deviation is $t$, the other fifteen deviations
sum to $-t$, so their squared sum is at least $t^2/15$.

The quadratic form $\mathbb E V=w^{\mathsf T}Bw$ is invariant under
permutations of the sixteen vertices. Its entries depend only on whether
two unordered edges are equal, share an endpoint, or are disjoint:
\[
 B_{ee}=\frac{19}{192},\qquad
 B_{ef}=\frac5{168}\quad(e\sim f),\qquad
 B_{ef}=-\frac{179}{17472}\quad(e\cap f=\varnothing).
\]
Here is a direct derivation of these rational counts. An edge whose least
common ancestor has $16,8,4,2$ leaves belongs to respectively
$16,8,4,2$ leaf paths; these four edge classes have sizes $64,32,16,8$.
Thus the trace of the leaf-centered covariance for a fixed tree is
\[
 32\frac14+16\frac3{16}+8\frac7{64}=\frac{95}{8}.
\]
Dividing by $120$ gives $19/192$. For any one leaf, the degrees of the
sixteen vertices in its path's edge set are $15,14,12,8$, with respective
multiplicities $2,2,4,8$. Each vertex's average degree over leaves is
$85/8$. The sum of these sixteen squared degree deviations is $495/4$.
Subtracting twice the covariance trace leaves the total covariance over
ordered adjacent edges, namely $495/4-2(95/8)=100$. Dividing by
$16\cdot15\cdot14=3360$ gives $5/168$. Finally, constant edge weights give
zero leaf variance. The ordered disjoint-edge covariance sum is therefore
$-95/8-100=-895/8$, and its $120\binom{14}{2}=10920$ terms give
$-179/17472$.

For an independent executable check, use the incidence indicators
$A_{\ell e}=\mathbf1_{e\in E_\ell}$ of one fixed tree, compute
\[
 \frac1{16}\sum_\ell A_{\ell e}A_{\ell f}
 -\left(\frac1{16}\sum_\ell A_{\ell e}\right)
  \left(\frac1{16}\sum_\ell A_{\ell f}\right),
\]
and average this value over the $120$ equal-edge cases, $1680$ unordered
adjacent-edge pairs, and $5460$ unordered disjoint-edge pairs. Relabeling
uniformly is exactly this orbit average. The supplied checker performs
all of these counts with rational arithmetic.

For a symbolic verification of the spectral bound, let $A$ be the adjacency
matrix of the line graph of $K_{16}$. Its eigenvalues are $28$ on constants,
$12$ on the fifteen-dimensional vertex component, and $-2$ on the
remaining $104$ dimensions. This also follows directly from the unsigned
vertex-edge incidence matrix $R$: $R^{\mathsf T}R=A+2I$ and
$RR^{\mathsf T}=14I+J$. Substituting the three rational entries above
therefore gives eigenvalues
\[
 0,\qquad \frac{33}{56},\qquad \frac{85}{2912}
\]
for $B$ on these respective subspaces. Consequently
\[
 \mathbb E V\le\frac{33}{56}\left(q-\frac{W^2}{120}\right).
\]
Jensen's inequality and then the two-dimensional Cauchy inequality give
\begin{align*}
 \min_{\text{trees}}\max_\ell L_\ell
 &\le\mathbb E\max_\ell L_\ell\\
 &\le\frac{17}{24}W+
       \sqrt{\frac{495}{56}}\sqrt{q-W^2/120}\\
 &\le\sqrt{120(17/24)^2+495/56}\sqrt q\\
 &\le\left(\sqrt{\frac{725}{1344}}+o(1)\right)K.
\end{align*}
This proves the lemma.
\end{proof}

\paragraph{An explicit deterministic choice.}
An exact subset dynamic program avoids enumerating tree relabelings.
Replace each $h_e$ by $b_e=\lceil\log_2N_e\rceil$, calculated from integer
bit lengths, and use unordered weights $w_{ab}=b_{ab}+b_{ba}$. This changes
any path weight by less than $170$ and hence preserves the coefficient.
For every subset $S$ of power-of-two cardinality let $F(S)$ be the minimum
possible maximum path weight for a balanced tree on $S$. For singleton sets
put $F(\{v\})=0$. If $W(S)$ is the sum of internal unordered weights, then
\[
 F(S)=\min_{\substack{A\sqcup B=S\\|A|=|B|}}
 \left[W(S)-W(A)-W(B)+\max\{F(A),F(B)\}\right].
\]
Indeed, the root cut weight is common to all paths, and the two subtree
problems optimize independently. Ancestor cuts, when present, are also
common additive constants and do not change this optimal substructure.
Store a minimizing split and recover a complete minimizing tree.

Compute $W(S)$ first for all $2^{16}$ masks. The dynamic program has
$\binom{16}{1}+\binom{16}{2}+\binom{16}{4}+\binom{16}{8}+1=14827$
states. Enumerating only splits with $\min S\in A$ examines
$120+1820\cdot3+12870\cdot35+6435=462465$ splits. All arithmetic is on
polynomial-bit integers. The selector is therefore deterministic and uses
polynomial time and space in $s$. This exact finite selector does not imply
a practical end-to-end SAT speedup. Once selected, the separator sum is
streamed depth first. The local recurrence is
\[
 T_d\le\operatorname{poly}(s)\,
 2^{(\beta_{\rm ps}+o(1))K_d}T_{d+1},\qquad
 \beta_{\rm ps}=\sqrt{725/1344}.
\]
All internal-node work is charged to the leaf evaluation count, and the
index arrays and accumulators remain polynomial in size.

\paragraph{Global schedule and coefficient.}
Retain the earlier polynomial-space schedule. Let
$\lambda=\lceil\log_2s\rceil$, $g=\lceil\sqrt\lambda\rceil$,
$B_0=\lfloor(\lambda-g)/80\rfloor$, and
\[
 u_0=B_0,\qquad u_d=\lfloor3B_0/4^d\rfloor\quad(d\ge1),\qquad
 a_d=\sum_{j=0}^d u_j+(d+1)\lceil6\log_2\lambda\rceil.
\]
Use leaf threshold $\lceil s/\lambda\rceil$. The total taper length is
at most $2B_0$. The symmetric root gate and the retained-side asymmetric
gate have a $\Theta(\sqrt\lambda)$ logarithmic margin. The fresh-side
and padding exponents are at most $9\lambda/40+o(\lambda)<\lambda$.
Near-lossless packing gives sharp child ratio $(1+o(1))/8$ and eventual
ratio at most $1/6$, so active depth is at most $\log_6\lambda+O(1)$
and $u_d\ge\lambda^{1-\log_6 4+o(1)}\to\infty$ there. The existing
buffer absorbs inherited precision exactly as in the prior proof.
Thus the unchanged schedule gives
\[
 \sum_d\frac{M_d}{\sqrt{u_d}}
 \le\left[\sqrt{80}\left(1+\frac1{3\sqrt3}\right)+o(1)\right]
 \frac{s}{\sqrt\lambda}.
\]
Multiplying by $(\pi\sqrt2/\ln2)\beta_{\rm ps}$ proves
Theorem~\ref{thm:adaptive-separator-new}. Leaves, tensor entry evaluation,
and polynomial factors have the same lower-order bounds as in the
inherited schedule. Every recursion is streamed, so space stays polynomial.

The new coefficient is the old $55.709478\ldots$ coefficient multiplied by
\[
 \frac{\sqrt{725/1344}}{\sqrt{170}/16}
 =\sqrt{\frac{290}{357}}=0.90129\ldots.
\]
This improves a displayed upper-bound certificate, not a measured running
time, and it leaves the SAT exponent's logarithmic order unchanged.
